%%
%% Conclusions
%%
\section{Conclusions}
\label{sec:conclusion}

% TODO: Completare

We have presented an extension to QuantumC that enables the compilation of C programs with vector and matrix operations into quantum circuits. Our approach introduces pattern-based recognition for common linear algebra operations and integrates dedicated intermediate representations into the existing MLIR-based pipeline.

The main contributions of this work include:
\begin{itemize}
    \item Pattern recognition for vector addition, dot product, and matrix multiplication
    \item VecMat and QAR intermediate representations for vector/matrix operations
    \item Support for hybrid programs combining scalar and vector/matrix code
    \item Compile-time array initialization
    \item Comparative analysis of QFT and ripple-carry arithmetic backends
\end{itemize}

\subsection{Limitations}
\label{sec:limitations}

The current implementation has several limitations:
\begin{itemize}
    \item Only three vector/matrix patterns are supported (vec\_add, vec\_dot, matmul)
    \item Array dimensions must be known at compile time
    \item The backend comparison uses structural circuit metrics at the logical gate level
\end{itemize}

\subsection{Future Work}
\label{sec:future}

Future directions include:
\begin{itemize}
    \item Supporting additional patterns (convolution, transpose, etc.)
    \item Circuit-level optimizations
    \item Integration with additional quantum backends
    \item Fault-tolerant analysis with gate decomposition
\end{itemize}
